\section{Background}

\subsection{Coverage Standards in Infrastructure Policy}

Infrastructure coverage standards establish minimum acceptable service levels for geographic access. Electricity coverage standards target 100\% of inhabited areas; telephone universal service requirements cover rural communities explicitly; the Rural Electrification Administration set a precedent in 1936 for using federal investment to close infrastructure gaps that market forces would not close \citep{REA1936}.

No equivalent coverage standard exists for highway access. The Federal Aid Highway Act of 1956 established the interstate system's scope, design standards, and funding mechanisms, but set no geographic coverage requirement \citep{FAHA1956}. The result was a system designed to connect existing urban centers rather than to achieve minimum access for the full population. The original route selection process, as documented by the Bureau of Public Roads, explicitly prioritized traffic volume over coverage \citep{BPR1947}.

The 30-mile standard we apply draws on published emergency response research \citep{Branas2005}, agricultural freight logistics literature \citep{USDA_FR2019}, and rural economic geography \citep{Partridge2008}. Emergency medicine research establishes that trauma patients transported more than 45 minutes from a trauma center have substantially worse outcomes; at rural road speeds of 50--55 mph, 45 minutes corresponds to approximately 40 miles, making 30 miles a conservative floor for effective emergency access. Agricultural freight research establishes that grain transport costs increase substantially beyond 25 miles from an elevator or major highway, making 30 miles a reasonable upper bound for cost-effective farm-to-highway access.

\subsection{Prior Gap Analyses}

The most comprehensive prior attempt to identify missing interstate links is the FHWA Future Interstate Study (2000), which identified 12,000 miles of potential new interstate corridors based on projected 2020 traffic volumes \citep{FHWA_FIS2000}. The study identified the I-69 Gulf-to-Midwest corridor, the Northern Tier corridor, and several Appalachian connectors as high-priority, but focused on volume-based justification rather than coverage.

The American Association of State Highway and Transportation Officials (AASHTO) maintains a list of proposed interstate designations \citep{AASHTO2023}, currently numbering approximately 50 proposed corridors. These designations reflect political as well as transportation rationale and have not been systematically ranked by coverage or economic impact.

State long-range transportation plans (LRTPs) identify state-level priority corridors but are not coordinated across state lines. A corridor like the Northern Tier, which would traverse Montana, North Dakota, Minnesota, Wisconsin, and Michigan, has separate state studies with no federal integration.

\subsection{Economic Geography of Highway Access}

The relationship between highway access and economic development is well-documented. \citet{Michaels2008} finds that US counties with interstate access have higher manufacturing employment and wages; \citet{Chandra2000} documents positive labor market effects from the interstate system's construction. More recently, \citet{Giroud2013} shows that plants in counties losing highway access (due to indirect routing) exhibit productivity declines.

The literature on rural economic geography consistently finds that isolation from the highway network compounds other disadvantages. \citet{USDA_ERS2020} documents that counties classified as ``deep rural'' (furthest from metro areas, typically also furthest from interstates) have the lowest per capita income growth, the highest poverty persistence, and the lowest business formation rates in the country. These are not coincidental correlations: distance from the highway network is a primary mechanism through which rural isolation transmits economic disadvantage.

Our C3 dimension (Economic Opportunity Access) operationalizes this relationship by measuring buffer GDP per capita relative to the national average. The finding that gap counties score 40\% higher on C3 than non-gap counties --- meaning their economic opportunity is substantially below the national average --- is consistent with this literature.

\subsection{Construction-Era Equity Harms}

Any analysis that frames new interstate investment as an equity benefit must acknowledge the equity harms of prior interstate construction. Between 1956 and 1980, the interstate system destroyed or severely damaged hundreds of urban neighborhoods, disproportionately affecting Black and low-income communities that lacked the political power to resist routing decisions. The Rondo neighborhood in St. Paul, Minnesota lost more than 600 homes and displaced hundreds of families to the construction of I-94 \citep{Ehrenfeucht2015}. Overtown in Miami, once known as the ``Harlem of the South,'' was effectively destroyed by the intersection of I-95 and I-395, which reduced the neighborhood's population from 40,000 to fewer than 10,000 \citep{Connolly2014}. The Tremé neighborhood in New Orleans was bisected by an elevated I-10 expressway over Claiborne Avenue, replacing a corridor of live oak trees and community gathering space \citep{Lewis1976}. These patterns were not aberrations. A 2021 analysis by the US Department of Transportation found that 73\% of interstate miles through urban areas were routed through communities of color \citep{USDOT_Equity2021}.

The Interstate 2.0 framework acknowledges this history and imposes an explicit alignment requirement for new corridors: proposed corridors that would route through existing communities must include community impact assessment, anti-displacement protections, and alignment alternatives before any federal designation is granted. This paper proposes corridors (Section~\ref{sec:corridors}) and ranks them by coverage efficiency; the alignment decisions for each corridor are outside this paper's scope but are a necessary second step before any construction commitment.

\subsection{Interstate 2.0 Context}

This paper is part of the ROUTE research module \citep{ROUTE_MODULE2026}, which scores 227 existing interstate corridors against a 16-dimensional rubric (v1.4) and uses gap analysis to identify investment priorities for Interstate 2.0. Paper A.1 \citep{ROUTE_A1} establishes the tier classification of existing corridors. This paper (B.1) identifies the gaps in that network. Papers B.2, B.3, and B.4 address bottlenecks, resilience holes, and intersection vulnerabilities respectively. The full Interstate 2.0 design framework appears in paper E.2 \citep{ROUTE_E2}.
